\documentclass{amsart}
% For dealing with references we use the comment environment
\usepackage{verbatim}
\newenvironment{reference}{\comment}{\endcomment}
%\newenvironment{reference}{}{}
\newenvironment{slogan}{\comment}{\endcomment}
\newenvironment{history}{\comment}{\endcomment}

% For commutative diagrams we use Xy-pic
\usepackage[all]{xy}

% We use 2cell for 2-commutative diagrams.
\xyoption{2cell}
\UseAllTwocells

% We use multicol for the list of chapters between chapters
\usepackage{multicol}

% This is generally recommended for better output
\usepackage{lmodern}
\usepackage[T1]{fontenc}

% For cross-file-references

% Package for hypertext links:
\usepackage{hyperref}

% For any local file, say "hello.tex" you want to link to please

% Theorem environments.
%
\theoremstyle{plain}
\newtheorem{theorem}[subsection]{Theorem}
\newtheorem{proposition}[subsection]{Proposition}
\newtheorem{lemma}[subsection]{Lemma}

\theoremstyle{definition}
\newtheorem{definition}[subsection]{Definition}
\newtheorem{example}[subsection]{Example}
\newtheorem{exercise}[subsection]{Exercise}
\newtheorem{situation}[subsection]{Situation}

\theoremstyle{remark}
\newtheorem{remark}[subsection]{Remark}
\newtheorem{remarks}[subsection]{Remarks}

\numberwithin{equation}{subsection}

% Macros
%
\def\lim{\mathop{\mathrm{lim}}\nolimits}
\def\colim{\mathop{\mathrm{colim}}\nolimits}
\def\Spec{\mathop{\mathrm{Spec}}}
\def\Hom{\mathop{\mathrm{Hom}}\nolimits}
\def\Ext{\mathop{\mathrm{Ext}}\nolimits}
\def\SheafHom{\mathop{\mathcal{H}\!\mathit{om}}\nolimits}
\def\SheafExt{\mathop{\mathcal{E}\!\mathit{xt}}\nolimits}
\def\Sch{\mathit{Sch}}
\def\Mor{\mathop{\mathrm{Mor}}\nolimits}
\def\Ob{\mathop{\mathrm{Ob}}\nolimits}
\def\Sh{\mathop{\mathit{Sh}}\nolimits}
\def\NL{\mathop{N\!L}\nolimits}
\def\CH{\mathop{\mathrm{CH}}\nolimits}
\def\proetale{{pro\text{-}\acute{e}tale}}
\def\etale{{\acute{e}tale}}
\def\QCoh{\mathit{QCoh}}
\def\Ker{\mathop{\mathrm{Ker}}}
\def\Im{\mathop{\mathrm{Im}}}
\def\Coker{\mathop{\mathrm{Coker}}}
\def\Coim{\mathop{\mathrm{Coim}}}

% Boxtimes
%
\DeclareMathSymbol{\boxtimes}{\mathbin}{AMSa}{"02}

%
% Macros for moduli stacks/spaces
%
\def\QCohstack{\mathcal{QC}\!\mathit{oh}}
\def\Cohstack{\mathcal{C}\!\mathit{oh}}
\def\Spacesstack{\mathcal{S}\!\mathit{paces}}
\def\Quotfunctor{\mathrm{Quot}}
\def\Hilbfunctor{\mathrm{Hilb}}
\def\Curvesstack{\mathcal{C}\!\mathit{urves}}
\def\Polarizedstack{\mathcal{P}\!\mathit{olarized}}
\def\Complexesstack{\mathcal{C}\!\mathit{omplexes}}
% \Pic is the operator that assigns to X its picard group, usage \Pic(X)
% \Picardstack_{X/B} denotes the Picard stack of X over B
% \Picardfunctor_{X/B} denotes the Picard functor of X over B
\def\Pic{\mathop{\mathrm{Pic}}\nolimits}
\def\Picardstack{\mathcal{P}\!\mathit{ic}}
\def\Picardfunctor{\mathrm{Pic}}
\def\Deformationcategory{\mathcal{D}\!\mathit{ef}}

\setlength{\emergencystretch}{3em}
\begin{document}
\title[Stacks Project, Tag 064Z]{351 --- Restriction of scalars detects pseudo-coherence}
\maketitle
\noindent Copyright (C) 2005--2025 Johan de Jong. Stacks Project authors. Source: \href{https://stacks.math.columbia.edu/tag/064Z}{Tag 064Z}.

\noindent Editorial difficulty note (not author text). Difficulty H2: The difficult direction builds a finite free approximation after extending scalars and lowers the highest nonzero cohomology degree by a cone. It proves that the cone preserves pseudo-coherence over the base and inductively recovers a finite approximation over the target algebra..

\noindent Editorial provenance note (not author text). History: excerpt from revision \texttt{a0444\allowbreak 6e57e\allowbreak c1fbc\allowbreak 252a8\allowbreak 71afc\allowbreak ec775\allowbreak 2fb28\allowbreak 07b14}, more-algebra.tex, lines 16282--16342. Reference presentation was adapted; original-author context and core support blocks were retained or added. The mathematical wording is unchanged. Licensed under GNU FDL 1.2 or later; no invariant sections or cover texts. See ../licenses/COPYING. Context blocks reproduce author definitions and supporting results; they are not additional counted problems.

\begin{lemma}
\label{lemma-finite-push-pseudo-coherent}
Let $A \to B$ be a ring map. Assume that $B$ is pseudo-coherent as an
$A$-module. Let $K^\bullet$ be a complex of $B$-modules.
The following are equivalent
\begin{enumerate}
\item $K^\bullet$ is $m$-pseudo-coherent
as a complex of $B$-modules, and
\item $K^\bullet$ is $m$-pseudo-coherent
as a complex of $A$-modules.
\end{enumerate}
The same equivalence holds for pseudo-coherence.
\end{lemma}

\begin{proof}
Assume (1). Choose a bounded complex of finite free $B$-modules
$E^\bullet$ and a map $\alpha : E^\bullet \to K^\bullet$ which is
an isomorphism on cohomology in degrees $> m$ and a surjection in degree $m$.
Consider the distinguished triangle
$(E^\bullet, K^\bullet, C(\alpha)^\bullet)$. By
Lemma \href{https://stacks.math.columbia.edu/tag/064W}{lemma recognize pseudo coherent (Tag 064W)}
$C(\alpha)^\bullet$ is $m$-pseudo-coherent as a complex of
$A$-modules. Hence it suffices to prove that $E^\bullet$ is
pseudo-coherent as a complex of $A$-modules, which follows from
Lemma \href{https://stacks.math.columbia.edu/tag/064Y}{lemma complex pseudo coherent modules (Tag 064Y)}.
The pseudo-coherent case of (1) $\Rightarrow$ (2) follows from this and
Lemma \href{https://stacks.math.columbia.edu/tag/064U}{lemma pseudo coherent (Tag 064U)}.

\medskip\noindent
Assume (2). Let $n$ be the largest integer such that $H^n(K^\bullet) \not = 0$.
We will prove that $K^\bullet$ is $m$-pseudo-coherent as a complex
of $B$-modules by induction on $n - m$. The case $n < m$ follows from
Lemma \href{https://stacks.math.columbia.edu/tag/064W}{lemma recognize pseudo coherent (Tag 064W)}.
Choose a bounded complex of finite free $A$-modules $E^\bullet$ and a
map $\alpha : E^\bullet \to K^\bullet$ which is an isomorphism on
cohomology in degrees $> m$ and a surjection in degree $m$.
Consider the induced map of complexes
$$
\alpha \otimes 1 : E^\bullet \otimes_A B \to K^\bullet.
$$
Note that $C(\alpha \otimes 1)^\bullet$ is acyclic in degrees
$\geq n$ as $H^n(E) \to H^n(E^\bullet \otimes_A B) \to H^n(K^\bullet)$
is surjective by construction and since $H^i(E^\bullet \otimes_A B) = 0$
for $i > n$ by the spectral sequence of
Example \href{https://stacks.math.columbia.edu/tag/0662}{example tor (Tag 0662)}.
On the other hand, $C(\alpha \otimes 1)^\bullet$
is $m$-pseudo-coherent as a complex of $A$-modules because
both $K^\bullet$ and $E^\bullet \otimes_A B$ (see
Lemma \href{https://stacks.math.columbia.edu/tag/064Y}{lemma complex pseudo coherent modules (Tag 064Y)})
are so, see
Lemma \href{https://stacks.math.columbia.edu/tag/064R}{lemma cone pseudo coherent (Tag 064R)}.
Hence by induction we see that $C(\alpha \otimes 1)^\bullet$
is $m$-pseudo-coherent as a complex of $B$-modules. Finally
another application of
Lemma \href{https://stacks.math.columbia.edu/tag/064R}{lemma cone pseudo coherent (Tag 064R)}
shows that $K^\bullet$ is $m$-pseudo-coherent as a complex of $B$-modules
(as clearly $E^\bullet \otimes_A B$ is pseudo-coherent as a complex
of $B$-modules). The pseudo-coherent case
of (2) $\Rightarrow$ (1) follows from this and
Lemma \href{https://stacks.math.columbia.edu/tag/064U}{lemma pseudo coherent (Tag 064U)}.
\end{proof}
\end{document}
